\Askhsh{13653}{2}
\begin{ylh}{1}
    \item 3.2. 1ο Κριτήριο ισότητας τριγώνων
    \item 3.7. Κύκλος - Μεσοκάθετος - Διχοτόμος
    \item 4.6. Άθροισμα γωνιών τριγώνου
    \item 5.9. Μια ιδιότητα του ορθογώνιου τριγώνου.
\end{ylh}
% ===================================================================
%  Αρχείο: 13653.tex (Fragment)
%  Αυτόματη μετατροπή μέσω doc2latex.py
% ===================================================================

ΘΕΜΑ 2

Σχεδιάζουμε γωνία $x\widehat{\text{O}}$y$\  = \ $60\textsuperscript{ο} και παίρνουμε σημείο Α επί της πλευράς Οx, τέτοιο ώστε AO$\  =$ 2. Φέρουμε τη διχοτόμο Οδ της γωνίας $x\widehat{\text{O}}$y και θεωρούμε σημείο Μ στην Οδ, τέτοιο ώστε AΜ$\  = \ $ΑΟ. Να υπολογίσετε:

\begin{alist}
\item Τη γωνία $\text{δ}\widehat{\text{O}}$y.

\monades{6}

\item Τις γωνίες του τριγώνου ΑΟΜ.

\monades{9}

\item Το μήκος του ύψους ΑB που αντιστοιχεί στη βάση ΟΜ του ισοσκελούς τριγώνου ΑΟΜ.

\monades{10}

\includesvg[width=4.58731in,height=3.93701in]{/home/spyros/Μαθηματικά/Trapeza_Thematwn/A_Lykeiou/Geometry/extracted/13653/media/image2.svg}
\end{alist}
